\chapter{Related Works}
\label{Chapter2}
% Table 1: Papers, Mechanics, and Methodology
\section{Traditional Methods}
\begin{table}[H]
	\centering
	\begin{tabular}{|l|p{12cm}|}
		\hline
		{\textbf{Paper}} & {\textbf{Mechanics}} \\ \hline
		Watershed &
		Flooding process on gradient topography with controlled basins \\ \hline
		Active Contour &
		Energy minimizing curve driven by internal smoothness and image forces \\ \hline
		GrabCut &
		Probabilistic graph cut with iterative foreground background modeling \\ \hline
		Region growing &
		Local similarity expansion based on pixel intensity or color distance \\ \hline
		BoxSup           & Weakly supervised semantic segmentation using bounding boxes and iterative region proposal refinement                 \\ \hline
		ScribbleSup      & Weak supervision via sparse scribble annotations propagated using graphical models \\ \hline
		Deep Extreme Cut & CNN predicts object mask from extreme point prompts                                \\ \hline
		Segment Anything & Uses a large image encoder based on Vision Transformers and a lightweight prompt encoder for points, boxes, or masks \\ \hline
	\end{tabular}
\end{table}

\begin{table}[H]
	\centering
	\begin{tabular}{|p{3cm}|p{4cm}|p{4cm}|p{4cm}|}
		\hline
		\multirow{2}{*}{\textbf{Paper}} &
		\multicolumn{3}{|c|}{\textbf{Methodology}} \\
		\cline{2-4}
		&
		\textbf{Input} &
		\textbf{Optimization} &
		\textbf{Final Labeling} \\ \hline
		Watershed &
		Image gradient, foreground and background markers &
		Flooding until basins meet &
		Watershed lines form region boundaries \\ \hline
		Active Contour &
		Image as continuous energy field, initial contour &
		Iterative energy minimization via curve evolution &
		Final curve represents object boundary \\ \hline
		GrabCut &
		Image as graph with color distributions (GMMs), bounding box &
		Repeated min cut optimization with model updates &
		Cut defines segmentation boundary \\ \hline
		Region growing &
		Image as pixel grid with feature vectors, seed pixels &
		Iterative neighbor expansion &
		Region boundary from grown region \\ \hline
		BoxSup &
		CNN feature maps extracted from region proposals and full images &
		Bounding boxes converted into soft foreground background supervision using objectness and low level segmentation cues &
		CNN trained on pseudo masks predicts dense semantic segmentation masks \\ \hline
		ScribbleSup &
		CNN feature maps over full image &
		Scribbles transformed into dense label maps via graph based or CRF propagation &
		CNN outputs semantic masks learned from propagated pseudo labels \\ \hline
		Deep Extreme Cut &
		CNN encodes image appearance along with spatial context &
		Extreme points encoded as distance maps concatenated with image channels &
		Single pass CNN inference produces object mask \\ \hline
		Segment Anything &
		Vision transformer generates dense image embeddings &
		Prompts encoded into tokens representing points, boxes, or masks &
		Transformer based mask decoder predicts multiple candidate masks with confidence scores \\ \hline
	\end{tabular}
\end{table}

\begin{table}[H]
	\centering
	\begin{tabular}{|p{3cm}|p{6cm}|p{6cm}|}
		\hline
		\textbf{Paper} & \textbf{Contribution} & \textbf{Limitation} \\ \hline
		Watershed &
		Introduced topographic interpretation of images and marker controlled segmentation to reduce over segmentation &
		Sensitive to noise and gradient quality \\ \hline
		Active Contour &
		Formalized segmentation as energy minimization over curves and influenced decades of variational methods &
		Struggles with concave boundaries and weak edges \\ \hline
		GrabCut &
		Demonstrated practical interactive segmentation with probabilistic modeling &
		Assumes color separability between foreground and background; fails when distributions overlap \\ \hline
		Region growing &
		Established seed based interactive segmentation &
		Highly sensitive to noise and intensity variation, leakage occurs at weak boundaries \\ \hline
		BoxSup & Leverages bounding boxes, which are approximately 15 times faster to collect than pixel-level masks & Bounding box contains no information about the specific shape of the object, causing "boundary error"-misalignment at the pixel level \\ \hline
		ScribbleSup & Scribbles are highly effective for annotating "stuff" categories that have no well-defined shape or boundary & The model is sensitive to the length and location of the scribbles \\ \hline
		Deep Extreme Cut & By using four points, DEXTR provides far more geometric information than a bounding box while remaining faster to annotate & Still struggles to capture extremely thin or "tiny" structures where the pixel area is minimal \\ \hline
		Segment Anything & Can handle ambiguous prompts by predicting multiple valid masks simultaneously, rather than averaging them into a single blurry result & The image encoder is a heavyweight Vision Transformer. Generating the initial image embedding is not yet "real-time" on lower-end devices \\ \hline
	\end{tabular}
\end{table}

\subsection{Watershed Transformation}
The Watershed transformation is a morphological segmentation tool based on a topographic analogy. It treats the gradient of an image as a relief surface where high intensities represent mountain ridges and low intensities represent valleys \cite{watershed_1979}.

\begin{enumerate}
	\item \textbf{Mechanics:} 
	The algorithm simulates the immersion of the topographic surface in water. Holes are pierced at local minima (or user markers), and as the surface is submerged, catchment basins fill up. At points where water from different basins would meet, "watershed lines" (dams) are constructed to partition the image.
	
	\item \textbf{Input:} 
	A grayscale image (typically the gradient magnitude) and a set of visual prompts (markers/points) identifying the foreground and background regions.
	
	\item \textbf{Output:} 
	A label map where each pixel is assigned to a specific marker's region, separated by one-pixel wide watershed lines.
\end{enumerate}

\subsection{Active Contour Models (Snakes)}
Active Contour Models, or "Snakes," are energy-minimizing splines that "slither" toward salient image features like edges and lines under the influence of internal and external forces. Snakes rely on a visual prompt—usually a rough initial contour drawn near the object of interest. The model then uses optimization to "shrink-wrap" itself onto the target's boundary \cite{snakes_1988}.

\begin{enumerate}
	\item \textbf{Mechanics:} 
	The snake's behavior is governed by an energy functional consisting of internal spline energy (maintaining smoothness and continuity) and external image energy (attracting the snake to edges). The goal is to find a contour position that minimizes the total Gibbs energy.
	
	\item \textbf{Input:} 
	The digital image and an initial set of points forming a rough contour (the visual prompt).
	
	\item \textbf{Output:} 
	A set of coordinates representing a smooth, optimized boundary that fits the object's edges.
\end{enumerate}


\subsection{GrabCut}
GrabCut is an interactive segmentation framework that extends Graph Cuts by using Gaussian Mixture Models (GMMs) and an iterative estimation procedure to extract foreground objects with minimal user effort. The primary visual prompt for GrabCut is a bounding box. The user drags a rectangle around the desired object; the area outside the box is marked as definite background, while the area inside is considered "unknown" and segmented automatically \cite{grab_cut_2004}.

\begin{enumerate}
	\item \textbf{Mechanics:} 
	The algorithm models the foreground and background color distributions using GMMs. It then constructs a Markov Random Field (MRF) where an energy function—balancing color similarity and spatial coherence—is minimized using a min-cut/max-flow algorithm.
	
	\item \textbf{Input:} 
	An RGB image and a bounding box (rectangle) prompt. Optional "brush strokes" can be added for refinement.
	
	\item \textbf{Output:} 
	A binary mask (foreground vs. background) and an alpha-matte for the object boundary.
\end{enumerate}

\subsection{Seeded Region Growing (SRG)}
Seeded Region Growing (SRG) is a pixel-level segmentation technique that partitions an image into regions based on the similarity of pixels to initial "seed" points. SRG is highly dependent on visual prompts. A user specifies one or more seed points for each region of interest. The algorithm then "grows" these regions by absorbing neighboring pixels that meet a similarity criterion \cite{region_growing_2014}.

\begin{enumerate}
	\item \textbf{Mechanics:} 
	The algorithm maintains a set of unallocated pixels that are neighbors to at least one grown region. At each step, the pixel with the smallest "distance" (highest similarity) to its adjacent region is added, and its neighbors are updated in the candidate list.
	
	\item \textbf{Output:} 
	A fully partitioned label map where each pixel is assigned to a region corresponding to one of the initial seeds.
\end{enumerate}


\section{Modern Methods}

\subsection{BoxSup}
BoxSup is a weakly-supervised framework designed to train semantic segmentation networks using bounding box annotations instead of expensive pixel-level masks \cite{box_sup_2015}.

\begin{enumerate}
	\item \textbf{Mechanics:} The method employs an iterative learning scheme. It uses unsupervised region proposal algorithms to generate candidate masks and then selects the most likely mask for each object based on the bounding box constraints and current network feedback.

	\item \textbf{Input:} Training images with bounding box annotations.
	
	\item \textbf{Output:} A trained Fully Convolutional Network (FCN) capable of producing dense semantic segmentation masks.
\end{enumerate}

\subsection{ScribbleSup}
ScribbleSup is a weakly-supervised approach that uses sparse scribbles to supervise the training of deep networks for semantic segmentation \cite{scribble_sup_2016}.

\begin{enumerate}
	\item \textbf{Mechanics:} The system uses a graphical model (CRF) built on super-pixels to propagate label information from the scribbles to the unknown, unmarked pixels. This propagation is guided by spatial constraints, appearance, and the network’s own semantic predictions.
	
	\item \textbf{Input:} RGB images and sparse scribble annotations (class-specific strokes).
	
	\item \textbf{Output:} A dense semantic label map for the entire image and a trained FCN model.
\end{enumerate}

\subsection{Deep Extreme Cut (DEXTR)}
Deep Extreme Cut (DEXTR) is an interactive segmentation method that converts four extreme points provided by a user into a precise object mask \cite{deep_extreme_cut_2018}.

\begin{enumerate}
	\item \textbf{Mechanics:} The four extreme points are encoded as a 2D Gaussian heatmap. This heatmap is concatenated with the original RGB image to form a 4-channel input, providing the CNN with an explicit spatial prior for the target object.
	
	\item \textbf{Input:} An RGB image and four extreme point clicks.
	
	\item \textbf{Output:} A precise, high-fidelity binary mask of the segmented object.
\end{enumerate}

\subsection{FocalClick}
FocalClick is a high-efficiency interactive segmentation framework designed for practical applications, particularly for mask correction and use on low-power devices \cite{focal_click_2022}.

\begin{enumerate}
	\item \textbf{Mechanics:} Upon receiving a user click, the model operates on two levels: a "Target Crop" (for coarse global segmentation) and a "Focus Crop" (a higher resolution local patch centered on the new click for fine refinement).
	
	\item \textbf{Input:} RGB image, positive/negative clicks, and an optional pre-existing (previous) mask.
	
	\item \textbf{Output:} An updated, refined segmentation mask that accurately reflects the new user clicks while maintaining global consistency.
\end{enumerate}

\bibliographystyle{IEEEtran} % Or IEEEtranN if using natbib
\bibliography{references}    % Replace 'references' with the name of your .bib file

%\begin{table}[htbp]
%	\caption{}
%	\begin{center}
%		\begin{tabular}{|c|c|c|c|}
%			\hline
%			\textbf{Table}&\multicolumn{3}{|c|}{\textbf{Table Column Head}} \\
%			\cline{2-4} 
%			\textbf{Head} & \textbf{\textit{Table column subhead}}& \textbf{\textit{Subhead}}& \textbf{\textit{Subhead}} \\
%			\hline
%			copy& More table copy$^{\mathrm{a}}$& &  \\
%			\hline
%			\multicolumn{4}{l}{$^{\mathrm{a}}$Sample of a Table footnote.}
%		\end{tabular}
%		\label{tab1}
%	\end{center}
%\end{table}




\begin{table}[htbp]
	\caption{Comparison Table}
	\begin{center}
		\begin{tabular}{|c|c|c|c|c|c|c|c|c|c|}
			\hline
			\multirow{2}{*}{\textbf{Papers}} & 
			\multirow{2}{*}{\textbf{Mechanics}} & \multicolumn{3}{|c|}{\textbf{Methodology}} & 
			\multicolumn{3}{|c|}{\textbf{Performance}} &
			\multirow{2}{*}{\textbf{Contribution}} & 
			\multirow{2}{*}{\textbf{Limitation}} \\
			\cline{3-8} & &
			\textbf{\textit{Stage 1}} & 
			\textbf{\textit{Stage 2}} & 
			\textbf{\textit{Stage 3}} & 
			\textbf{\textit{Data}} & 
			\textbf{\textit{Metric}} &
			\textbf{\textit{Result}} &
			&  \\
			
			
			\hline
			 copy& More table copy$^{\mathrm{a}}$& & & & & & & & \\
			\hline
			\multicolumn{4}{l}{$^{\mathrm{a}}$Sample of a Table footnote.}
		\end{tabular}
		\label{tab1}
	\end{center}
\end{table}
